\vspace{2.5pt plus 3pt minus 0.5pt}\noindent\hspace*{1.2em}\kbvnum{(3)}\hspace{0.55em}ಎರಡು ಪದಗಳುಳ್ಳ  ಐದು ವಾಕ್ಯಗಳನ್ನು ಜ್ಞಾಪಿಸಿರಿ.\par

\noindent\hspace*{1.2em}\kbvlabel{ಉದಾ}\hspace{0.6em}ಹೂವು ಅರಳಿತು...\par

\vspace{2.5pt plus 3pt minus 0.5pt}\noindent\hspace*{1.2em}\kbvnum{(4)}\hspace{0.55em}ಮೂರು ಪದಗಳುಳ್ಳ ಐದು ವಾಕ್ಯಗಳನ್ನು ಬರೆಯಿರಿ.\par

\noindent\hspace*{1.2em}\kbvlabel{ಉದಾ}\hspace{0.6em}ಬೆಕ್ಕು ಹಾಲು ಕುಡಿಯಿತು...\par

\vspace{2.5pt plus 3pt minus 0.5pt}\noindent\hspace*{1.2em}\kbvnum{(5)}\hspace{0.55em}ನಾಲ್ಕು ಪದಗಳುಳ್ಳ ಮೂರು ವಾಕ್ಯಗಳನ್ನು ಬರೆಯಿರಿ.\par

\vspace{2.5pt plus 3pt minus 0.5pt}\noindent\hspace*{1.2em}\kbvnum{(6)}\hspace{0.55em}ಈ ಪದಗಳನ್ನು ಸ್ಪಷ್ಟವಾಗಿ ಉಚ್ಚರಿಸಿರಿ : ಹಕ್ಕಿ, ಅಕ್ಕಿ, ಅನ್ನ, ಅಣ್ಣ, ಉಷ್ಣ, ತೀಕ್ಷ್ಣ, ಗಾನ, ಜ್ಞಾನ, ರಾಷ್ಟ್ರ, ಆಜ್ಞೆ, ನಕ್ಷತ್ರ, ಉತ್ಥಾನ, ಉತ್ಸವ, ಸಂಸ್ಕೃತ, ಅಷ್ಟಾವಕ್ರ.\par

\vspace{2.5pt plus 3pt minus 0.5pt}\noindent\hspace*{1.2em}\kbvnum{(7)}\hspace{0.55em}ದಿಹಳ, ರಸಳ, ನಿಜರ, ಜನನ, ವತಾರೆ, ಮರಿಪಳ, ನಿವಾಶರ, ಪಾಜಾತರಿ, ಯಲವಾಶಿ, ಮಣರಕವೆಂಟ  - ಈ ಪದಗಳಲ್ಲಿ ಅಕ್ಷರಗಳು ಕಲಸಿಹೋಗಿವೆ ; ಅರ್ಥವತ್ತಾದ ಪದಗಳಾಗುವಂತೆ ಅಕ್ಷರಗಳನ್ನು ಕ್ರಮವಾಗಿ ಜೋಡಿಸಿ ಬರೆಯಿರಿ.\par

\vspace{2.5pt plus 3pt minus 0.5pt}\noindent\hspace*{1.2em}\kbvnum{(8)}\hspace{0.55em}ಅ-ಪಟ್ಟಿಯಲ್ಲಿರುವ ಪದಗಳೊಂದಿಗೆ ಬ-ಪಟ್ಟಿಯಲ್ಲಿರುವ ಪದಗಳನ್ನು ಕೂಡಿಸಿ ಅರ್ಥಪೂರ್ಣವಾದ ವಾಕ್ಯಗಳನ್ನು ಬರೆಯಿರಿ.\par

\vspace{5pt plus 4pt minus 1pt}
\begin{center}
{\renewcommand{\arraystretch}{1.24}\setlength{\tabcolsep}{10pt}
\begin{tabular}{l c l}
\textbf{ಅ-ಪಟ್ಟಿ} & \hspace{12mm} & \textbf{ಬ-ಪಟ್ಟಿ} \\[3pt]
ಕೋಳಿ & \ldots\ldots & ಬೊಗಳಿತು \\
ಕತ್ತೆ & \ldots\ldots & ಘರ್ಜಿಸಿತು \\
ನಾಯಿ & \ldots\ldots & ಅರಚಿತು \\
ಕುದುರೆ & \ldots\ldots & ಕೂಗಿತು \\
ಸಿಂಹ & \ldots\ldots & ಕೆನೆಯಿತು \\
\end{tabular}}
\end{center}
\vspace{5pt plus 4pt minus 1pt}

\vspace{2.5pt plus 3pt minus 0.5pt}\noindent\hspace*{1.2em}\kbvnum{(9)}\hspace{0.55em}ಕೆಳಗೆ ಕೊಟ್ಟಿರುವ ವಾಕ್ಯಗಳಲ್ಲಿ ಪದಗಳು ಸ್ಥಾನಪಲ್ಲಟಗೊಂಡಿವೆ. ಅವುಗಳನ್ನು ಕ್ರಮವಾಗಿ ಜೋಡಿಸಿ ಬರೆಯಿರಿ.\par

\vspace{2.5pt plus 3pt minus 0.5pt}\noindent\hspace*{1.2em}\kbvnum{(೧)}\hspace{0.55em}ಜೋಳದ ಆಹಾರ ರುಚಿಯಾದ ರೊಟ್ಟಿ.\par

\vspace{2.5pt plus 3pt minus 0.5pt}\noindent\hspace*{1.2em}\kbvnum{(೨)}\hspace{0.55em}ಅತ್ಯುತ್ತಮವಾದ ಮಗುವಿಗೆ ತಾಯಿಯ ಹಾಲು ಆಹಾರ.\par

\vspace{2.5pt plus 3pt minus 0.5pt}\noindent\hspace*{1.2em}\kbvnum{(3)}\hspace{0.55em}ಕಡ್ಡಾಯವಾಗಿ ಕರ್ನಾಟಕದಲ್ಲಿರುವ ಕನ್ನಡವನ್ನು ಕನ್ನಡೇತರರು ಕಲಿಯಬೇಕು.\par
